\documentclass[10pt,a4paper]{amsart}
\usepackage[margin=19mm]{geometry}
\usepackage{amsmath,amssymb,mathtools}
\usepackage[scr=rsfs,cal=euler]{mathalfa}
\usepackage{xfrac}
\usepackage[unicode,hidelinks,bookmarks=false]{hyperref}
\newcommand{\ie}{i.e.\ }
\newcommand{\cf}{cf.\ }
\newcommand{\resp}{respectively\ }
\newcommand{\Subsection}{\subsection}
\providecommand{\nobreakdash}{\nobreak\hbox{-}}
\setlength{\emergencystretch}{2em}
\multlinegap0pt
\usepackage{mathtools}
\usepackage{amssymb}
\usepackage{tikz}
\usepackage[shortlabels]{enumitem}
\usepackage{float}
\usepackage{xcolor}
\newtheorem{lemma}{Lemma}
\newtheorem{proposition}{Proposition}
\usepackage{cleveref}
\crefname{lemma}{Lemma}{Lemmas}
\crefname{proposition}{Proposition}{Propositions}
\newcommand{\Cart}{\cat{C}^{\infty}\cat{Cart}}
\newcommand{\CsSet}{\textnormal{C}^{\infty}\sset}
\newcommand{\cat}[1]{\mathsf{#1}}
\newcommand{\femb}[1]{\mathfrak{C}^{\infty}\mathfrak{Emb}_{#1}}
\newcommand{\fembnoniso}[1]{\cat{C}^{\infty}\cat{Emb}_{#1}}
\newcommand{\fieldnoniso}[1]{\cat{Field}_{#1}}
\newcommand{\id}{\mathrm{id}}
\newcommand{\psh}{\mathrm{PSh}}
\newcommand{\sset}{\cat{sSet}}
\title{Moduli spaces of geometric functorial field theories}
\author{Jacek Kenig, Dmitri Pavlov}
\date{Source: arXiv:2609.11812v1 (CC BY 4.0)}
\begin{document}
\maketitle

\noindent\emph{Source and licence.} Moduli spaces of geometric functorial field theories. Version arXiv:2609.11812v1 (CC BY 4.0), distributed by arXiv under the Creative Commons Attribution 4.0 International licence, http://creativecommons.org/licenses/by/4.0/, as recorded in the arXiv licence metadata for this exact version and shown on the abstract page https://arxiv.org/abs/2609.11812v1. Full licence notice: \texttt{licenses/arxiv-2609.11812-CC-BY-4.0.txt}.\par\smallskip

\noindent\textit{Editorial scope.} Counted unit: the article's proposition prop:reduced.category.discrete (source0.tex lines 3051--3079). The article's complete original proof of it is delivered with it (source0.tex lines 3081--3324, one \texttt{proof} environment of 244 lines). No mathematics is written, repaired or re-derived: the statement and the proof are the article's own lines, in the article's own notation, and no retained line is substituted or re-flowed.  Retained uncounted as the source-local context the proof needs: lines 2951--2971.  Nothing was omitted from any retained span, so the package carries no omission record.  No retained line contains a citation, so the wrapper carries no bibliography and no bibliography entry.  No retained line contains a figure environment, an image-inclusion command or drawing code, so no image is delivered and the package declares no asset dependency.  Editorial formatting only, in addition: an AMS \texttt{amsart} preamble that restates the article's own declarations and macros used by the retained lines, namely the environments \texttt{lemma}, \texttt{proposition} on separate counters, together with the standard packages those lines need.  The retained environments print in this document as Lemma 1, Proposition 1 in order of appearance, whereas the article prints the same items with the numbering of its own full document.  The complete native source of the article as distributed in the arXiv e-print for this exact version is bundled unchanged as \texttt{sources/210051/source0.tex}.
\begin{lemma}\label{lem:LKE.raw.triples}
  Let \(F\in \fieldnoniso{d}\) be a geometric structure without isotopies.
  Let \(q\in \femb{d}\), \(L\in \Cart\), and \(k\geq 0\). Then
  \[
    \bigl(\mathfrak{I}_{d}F(q)\bigr)_{L,k}
  \]
  is naturally identified with the quotient of
  \[
    \coprod_{p\in \fembnoniso{d}}
    F(p)_k\times \femb{d}(q,\iota(p))_L
  \]
  by the equivalence relation generated by
  \[
    (p_1,F(h)(x),e)
    \sim
    (p_2,x,(\iota(h)\times \id_L)\circ e)
  \]
  for every morphism \(h:p_1\to p_2\) in \(\fembnoniso{d}\),
  \(x\in F(p_2)_k\), and
  \(e\in \femb{d}(q,\iota(p_1))_L\).
\end{lemma}

\begin{proposition}\label{prop:reduced.category.discrete}
  Let \(F\in \fieldnoniso{d}\), \(q\in \femb{d}\),
  \(L\in \Cart\), and \(k\geq 0\). Then the set of
  \((L,k)\)-simplices of \(\mathfrak{I}_{d}F(q)\) is naturally identified
  with the set of isomorphism classes of triples
  \[
    (p,x,e),
  \]
  where \(p\in \fembnoniso{d}\), \(x\in F(p)_k\), and
  \[
    e\in \femb{d}(q,\iota(p))_L
  \]
  The map \(e\) is a base-preserving \(L\)-parameterized family of
  fiberwise open embeddings whose joint image is \(p\).

  Two such triples \((p,x,e)\) and \((p',x',e')\) are isomorphic if and only
  if there exists a necessarily unique isomorphism \(h:p\to p'\) in
  \(\fembnoniso{d}\) such that
  \[
    F(h)(x')=x,
    \qquad
    e'=(\iota(h)\times \id_L)\circ e.
  \]
  More generally, the same description of connected components by
  isomorphism classes of reduced triples holds with \(F(p)_k\) replaced by
  \(X(p)_{L,k}\), for every
  \(X\in\psh(\fembnoniso{d},\CsSet)\). The reduction depends only on the
  embedding datum \(e\).
\end{proposition}

\begin{proof}
  Fix \(q\in \femb{d}\), \(L\in \Cart\), \(k\geq 0\), and
  \(X\in \psh(\fembnoniso{d},\CsSet)\). Let
  \(\mathcal{G}_{X,q,L,k}\) be the category whose objects are triples
  \((p,x,e)\) with
  \[
    p\in \fembnoniso{d},\qquad
    x\in X(p)_{L,k},\qquad
    e\in \femb{d}(q,\iota(p))_L,
  \]
  Its morphisms
  \[
    (p_1,x_1,e_1)\longrightarrow (p_2,x_2,e_2)
  \]
  are morphisms \(h:p_1\to p_2\) in \(\fembnoniso{d}\) satisfying
  \[
    X(h)(x_2)=x_1,
    \qquad
    e_2=(\iota(h)\times \id_L)\circ e_1.
  \]
  Let \(\mathcal{G}^{\mathrm{red}}_{X,q,L,k}\) be the full subcategory on
  the triples for which \(e\) is base-preserving and has joint image equal
  to \(p\). We first prove that its inclusion into
  \(\mathcal{G}_{X,q,L,k}\) induces a bijection on connected components and
  that \(\mathcal{G}^{\mathrm{red}}_{X,q,L,k}\) is a groupoid with at most
  one morphism between any two objects.

  First, every representative may be chosen so that the map~\(e\) is
  base-preserving. Let \((p,x,e)\) be an object of
  \(\mathcal{G}_{X,q,L,k}\), and write
  \[
    e=(f,g\times \id_L):q\times \id_L\longrightarrow p\times \id_L.
  \]
  Since morphisms in \(\femb{d}\) are fiberwise open embeddings over maps of
  base spaces, the map \(e\) canonically factors through the pullback of \(p\)
  along the base map \(g\times \id_L\). In other words, if
  \(p'\coloneqq g^{*}p\), then there is a factorization
  \[
    q\times \id_L
    \xrightarrow{\ \widetilde{e}\ }
    p'\times \id_L
    \xrightarrow{\ \iota(h)\times \id_L\ }
    p\times \id_L,
  \]
  where \(\widetilde{e}\) lies over the identity on the base and
  \(h:p'\to p\) is the canonical morphism. Passing from \((p,x,e)\) to
  \((p',X(h)(x),\widetilde{e})\) gives a morphism
  \[
    (p',X(h)(x),\widetilde{e})\longrightarrow (p,x,e)
  \]
  in \(\mathcal{G}_{X,q,L,k}\). Thus every connected component contains an
  object whose embedding family is base-preserving.

  Second, after making the representative base-preserving, one may further
  replace the target by the joint image of the family of embeddings. Suppose
  therefore that the map
  \[
    e:q\times \id_L\longrightarrow p\times \id_L
  \]
  lies over the identity on the base. For each \(\ell\in L\), let
  \(e_\ell:q\to p\) denote the corresponding fiberwise open embedding. Define
  \[
    p^\circ\coloneqq \operatorname{im}(\pi_1\circ e)\subset p,
  \]
  where \(\pi_1:p\times \id_L\to p\) is the projection. Equivalently,
  \[
    p^\circ=\bigcup_{\ell\in L}\operatorname{im}(e_\ell).
  \]
  Since each \(e_\ell\) is an open embedding, \(p^\circ\) is an open subobject
  of \(p\), hence again an object of \(\fembnoniso{d}\). The map \(e\) factors
  through \(p^\circ\times \id_L\) as
  \[
    q\times \id_L
    \xrightarrow{\ e^\circ\ }
    p^\circ\times \id_L
    \xrightarrow{\ \iota(j)\times \id_L\ }
    p\times \id_L,
  \]
  where \(j:p^\circ\hookrightarrow p\) is the open inclusion. Passing from
  \((p,x,e)\) to \((p^\circ,X(j)(x),e^\circ)\) gives a morphism
  \[
    (p^\circ,X(j)(x),e^\circ)\longrightarrow (p,x,e)
  \]
  in \(\mathcal{G}_{X,q,L,k}\). By construction, the new representative has
  the property that its target is exactly the union of the images of the
  family \(\{e_\ell\}_{\ell\in L}\).

  Applying the two reductions successively defines a functor
  \[
    \nu:\mathcal{G}_{X,q,L,k}\longrightarrow
    \mathcal{G}^{\mathrm{red}}_{X,q,L,k}.
  \]
  On objects, \(\nu\) sends a triple \((p,x,e)\) to its reduced representative
  \[
    (\overline{p},\overline{x},\overline{e}),
    \qquad
    \overline{x}\coloneqq X(r)(x),
  \]
  where \(r:\overline{p}\to p\) is the composite reduction morphism.

  We now define \(\nu\) on morphisms. Let
  \[
    a:(p_1,x_1,e_1)\longrightarrow(p_2,x_2,e_2)
  \]
  be a morphism in \(\mathcal{G}_{X,q,L,k}\). For \(i=1,2\), write
  \[
    r_i:\overline{p}_i\longrightarrow p_i
  \]
  for the composite reduction morphism, and write
  \[
    (\overline{p}_i,\overline{x}_i,\overline{e}_i)
  \]
  for the reduced representative of \((p_i,x_i,e_i)\). Thus
  \[
    \overline{x}_i=X(r_i)(x_i),
    \qquad
    e_i=(\iota(r_i)\times\id_L)\circ \overline{e}_i.
  \]

  The identity
  \[
    e_2=(\iota(a)\times\id_L)\circ e_1
  \]
  implies that \(\iota(a)\times\id_L\) sends the joint image of
  \(\overline{e}_1\) into the joint image of \(\overline{e}_2\). Since
  \(\overline{p}_i\) is defined as this joint image after making the embedding
  base-preserving, there is a unique induced morphism
  \[
    \overline{a}:\overline{p}_1\longrightarrow\overline{p}_2
  \]
  such that
  \[
    r_2\overline{a}=a r_1.
  \]

  It remains to check that \(\overline{a}\) is a morphism of triples, i.e., that
  it
  satisfies the condition on the \(X\)-coordinates. Since \(X\) is
  contravariant, we have
  \[
    X(\overline{a})(\overline{x}_2)
    =
    X(\overline{a})\bigl(X(r_2)(x_2)\bigr)
    =
    X(r_1)\bigl(X(a)(x_2)\bigr)
    =
    X(r_1)(x_1)
    =
    \overline{x}_1.
  \]
  Hence \(\overline{a}\) defines a morphism
  \[
    (\overline{p}_1,\overline{x}_1,\overline{e}_1)
    \longrightarrow
    (\overline{p}_2,\overline{x}_2,\overline{e}_2)
  \]
  in \(\mathcal{G}^{\mathrm{red}}_{X,q,L,k}\). The uniqueness of the induced
  morphism \(\overline{a}\) makes this assignment
  compatible with identities and composition, so \(\nu\) is a functor.

  If \(J:\mathcal{G}^{\mathrm{red}}_{X,q,L,k}\hookrightarrow
  \mathcal{G}_{X,q,L,k}\) denotes the inclusion, then \(\nu J=\id\). Moreover,
  the morphisms produced by the two reduction steps give a natural
  transformation
  \[
    J\nu\longrightarrow \id_{\mathcal{G}_{X,q,L,k}}.
  \]
  Therefore \(J\) induces a bijection on connected components: every object of
  \(\mathcal{G}_{X,q,L,k}\) is connected to a reduced object, and any zigzag
  between reduced objects may be sent by \(\nu\) to a zigzag entirely inside
  \(\mathcal{G}^{\mathrm{red}}_{X,q,L,k}\).

  It remains to prove the stated property of the reduced category. Let
  \[
    h:(p_1,x_1,e_1)\longrightarrow (p_2,x_2,e_2)
  \]
  be a morphism in \(\mathcal{G}^{\mathrm{red}}_{X,q,L,k}\). Since both
  embedding families are base-preserving, the morphism \(h:p_1\to p_2\) lies
  over the identity on the common base. Since
  \[
    e_2=(\iota(h)\times \id_L)\circ e_1
  \]
  and \(e_2\) is jointly surjective, every point of \(p_2\) lies in the image
  of \(h\). Thus \(h\) is surjective on total spaces. But \(h\) is already a
  fiberwise open embedding over the fixed base, so it is an isomorphism.
  Therefore every morphism in
  \(\mathcal{G}^{\mathrm{red}}_{X,q,L,k}\) is an isomorphism.

  We now show that between any two objects there is at most one morphism. Let
  \[
    h_1,h_2:(p_1,x_1,e_1)\longrightarrow (p_2,x_2,e_2)
  \]
  be two morphisms in \(\mathcal{G}^{\mathrm{red}}_{X,q,L,k}\). By definition,
  \[
    e_2=(\iota(h_1)\times \id_L)\circ e_1
    =(\iota(h_2)\times \id_L)\circ e_1.
  \]
  Let \(y_1\in p_1\). Since \(e_1\) is jointly surjective, there exist
  \(\ell\in L\) and \(y\in q\) such that \(e_{1,\ell}(y)=y_1\). Then
  \[
    h_1(y_1)
    =h_1(e_{1,\ell}(y))
    =e_{2,\ell}(y)
    =h_2(e_{1,\ell}(y))
    =h_2(y_1).
  \]
  Since this holds for every \(y_1\in p_1\), we conclude that \(h_1=h_2\).

  In particular, every automorphism is the identity. Since
  \(\mathcal{G}^{\mathrm{red}}_{X,q,L,k}\) is a groupoid with at most one
  morphism between any two objects, it is equivalent to the discrete category
  of its isomorphism classes. Explicitly, if
  \[
    \pi_0\bigl(\mathcal{G}^{\mathrm{red}}_{X,q,L,k}\bigr)
  \]
  denotes the set of isomorphism classes of objects, then the quotient functor
  \[
    Q:
    \mathcal{G}^{\mathrm{red}}_{X,q,L,k}\longrightarrow
    \pi_0\bigl(\mathcal{G}^{\mathrm{red}}_{X,q,L,k}\bigr)
  \]
  is essentially surjective by construction, and it is fully faithful because
  for any two objects there is exactly one morphism between them if they are
  isomorphic, and none otherwise.
  Now take \(X=\lambda^\infty F\). By
  \Cref{lem:LKE.raw.triples},
  \(\bigl(\mathfrak{I}_{d}F(q)\bigr)_{L,k}\) is the set of connected
  components of \(\mathcal{G}_{\lambda^\infty F,q,L,k}\), and
  \[
    (\lambda^\infty F)(p)_{L,k}=F(p)_k.
  \]
  The preceding argument therefore identifies this set with the isomorphism
  classes of the triples in the statement. The description of an isomorphism
  between two such triples follows directly from the definition of the
  morphisms in \(\mathcal{G}_{\lambda^\infty F,q,L,k}\) and their uniqueness
  in the reduced category.

  Finally, both reductions are defined using only \(e\): the first replaces
  the target by the pullback determined by the base map of \(e\), and the
  second replaces that pullback by the joint image of the resulting
  base-preserving family. The \(x\)-coordinate is only transported along the
  resulting morphism. Hence the reduction is independent of \(x\) and of
  \(k\), and the same argument applies to arbitrary \(X\).
\end{proof}

\end{document}
